\section{稳定性、误差传播与证书层级}

\subsection{误差源必须分层}
最终输出误差不能用单一 neural validation loss 代表。生产系统至少包含
\begin{equation}
\boxed{
E_{\rm total}
\lesssim
E_{\rm model}
+E_{\rm domain}
+E_{\rm discretization}
+E_{\rm solve}
+E_{\rm tensor\,surrogate}
+E_{\rm thermal\,ROM}
+E_{\rm time}.
}
\end{equation}
其中：
\begin{itemize}
  \item $E_{\rm model}$：filament conductor、材料 constitutive law、热边界等物理模型误差；
  \item $E_{\rm domain}$：有限电磁/热计算域和边界截断误差；
  \item $E_{\rm discretization}$：Maxwell/thermal 空间离散与几何 deposition 误差；
  \item $E_{\rm solve}$：truth linear/nonlinear solve 的代数误差；
  \item $E_{\rm tensor\,surrogate}$：$Z/H$ 参数代理误差；
  \item $E_{\rm thermal\,ROM}$：thermal projection/truncation error；
  \item $E_{\rm time}$：ETD/IMEX 时间离散误差。
\end{itemize}
不同量纲和范数的误差不能直接数值相加；上式表示经相应稳定性/敏感度常数传播后的结构化预算。

\subsection{Joule tensor 误差到热源误差}
令
\begin{equation}
\Delta H_j=\widehat H_j-H_j.
\end{equation}
则任意端口电流 $c$ 下
\begin{equation}
\left|\widehat q_{{\rm vol},j}-q_{{\rm vol},j}\right|
=\left|\operatorname{Re}(c^H\Delta H_jc)\right|
\le \|\Delta H_j\|_2\,\|c\|_2^2.
\end{equation}
因此只要 operating domain 给出 $\|c\|\le C_c$，tensor operator-norm error 就能直接给出整个 current space 的热源误差上界，不需要枚举所有电流组合。

\subsection{阻抗代理误差}
对
\begin{equation}
\Delta Z_{\rm field}=\widehat Z_{\rm field}-Z_{\rm field},
\end{equation}
有
\begin{equation}
\|\Delta v_{\rm field}\|_2
\le\|\Delta Z_{\rm field}\|_2\|c\|_2,
\end{equation}
以及平均功率误差
\begin{equation}
\left|\Delta P_{\rm field}\right|
\le\frac12\|\Delta Z_{\rm field}\|_2\|c\|_2^2.
\end{equation}
线圈电阻误差必须单独归入 constitutive/model error，不与 $Z_{\rm field}$ neural error 混合。

\subsection{热动力学误差传播}
真实 reduced model 写为
\begin{equation}
M\dot a=-Ka+q_{\rm vol}(g,c)+q_{\rm wire}(a,g,c)+f_T,
\end{equation}
代理模型只把 $q_{\rm vol}$ 换成 $\widehat q_{\rm vol}$。令 $e=\hat a-a$。若目标 thermal state domain 内
\begin{equation}
F(a)=M^{-1}\left[-Ka+q_{\rm wire}(a,g,c)+f_T\right]
\end{equation}
具有 Lipschitz constant $L_F$，并满足
\begin{equation}
\|M^{-1}(\widehat q_{\rm vol}-q_{\rm vol})\|\le\varepsilon_F,
\end{equation}
则 $e(0)=0$ 时
\begin{equation}
\|e(t)\|
\le
\begin{cases}
\varepsilon_F\dfrac{e^{L_Ft}-1}{L_F},&L_F>0,\\[0.7em]
\varepsilon_F t,&L_F=0.
\end{cases}
\end{equation}
该界不包含 thermal ROM/time discretization error；这些必须另行加入误差预算。

\subsection{收缩情形}
若真实 reduced thermal dynamics 在目标域某能量范数下满足
\begin{equation}
\mu(J_F)\le-\gamma<0,
\end{equation}
则
\begin{equation}
\|e(t)\|
\le\frac{\varepsilon_F}{\gamma}
\left(1-e^{-\gamma t}\right)
\le\frac{\varepsilon_F}{\gamma}.
\end{equation}
只有在真实 wire-heating feedback 与 thermal operators 上验证该条件后，才允许宣称长期误差饱和；不能因为扩散项本身耗散就默认整个电热系统 contractive。

\subsection{thermal ROM error}
thermal basis 的主要数值证据必须来自实际 Galerkin residual：
\begin{equation}
r_K=Q-K\Phi(\Phi^TK\Phi)^{-1}\Phi^TQ,
\end{equation}
以及相应 dynamic residual。冻结 test geometries 上还应直接比较 full thermal 与 reduced thermal 的温度、最大温升和稳态结果。

\subsection{输出误差}
温度场满足
\begin{equation}
T=T_{\rm amb}+\Phi a+T_{\rm tail}.
\end{equation}
因此
\begin{equation}
\|\widehat T-T\|_\infty
\le
\|\Phi\|_{2\to\infty}\|e\|_2
+\|T_{\rm tail}\|_\infty
+E_{\rm time,T}.
\end{equation}
阻抗误差则由 $Z_{\rm field}$ surrogate 与 $R_{\rm wire}$ constitutive error 分别控制。

\subsection{证书层级}
所有 ``certificate'' 或 ``validated'' 字样必须区分：
\begin{enumerate}
  \item \textbf{严格证书：}数学上覆盖完整声明参数域且前提已验证；
  \item \textbf{条件定理：}例如上述 contraction bound，只在条件成立时生效；
  \item \textbf{冻结测试验证：}有限 geometry/current/state test set 上的最大误差；
  \item \textbf{经验诊断：}POD 谱、训练/验证 loss、sampling envelope 等。
\end{enumerate}
Physics Gate 的数值收敛实验通常属于冻结验证或数值证据，不得自动升级为全域严格证书。
